\documentclass[11pt,a4paper,fontset=fandol]{ctexart}
\usepackage[a4paper,top=2.3cm,bottom=2.5cm,left=2.3cm,right=2.3cm,headheight=18pt,headsep=0.55cm,footskip=25pt]{geometry}
\usepackage{amsmath,amssymb,amsthm,fancyhdr,lastpage,enumitem,titlesec,xcolor,hyperref,bookmark}
\usepackage[most]{tcolorbox}
\usepackage{setspace,array}
\hypersetup{colorlinks=true,linkcolor=blue!50!black,urlcolor=blue!50!black}
\setstretch{1.15}
\setlength{\parindent}{2em}
\setlength{\parskip}{0.28em}
\allowdisplaybreaks
\setlist[enumerate]{itemsep=0.42em,topsep=0.35em,leftmargin=2.35em}
\setlist[itemize]{itemsep=0.25em,topsep=0.25em,leftmargin=1.9em}
\definecolor{mainblue}{RGB}{36,83,140}
\definecolor{lightblue}{RGB}{241,246,252}
\definecolor{lightgray}{RGB}{246,246,246}
\definecolor{accent}{RGB}{153,76,0}
\definecolor{rulegray}{RGB}{110,118,128}
\titleformat{\section}{\Large\bfseries\color{mainblue}}{\thesection.}{0.45em}{}[\vspace{0.2em}\color{rulegray}\titlerule]
\titlespacing*{\section}{0pt}{1.15em}{0.75em}
\titleformat{\subsection}{\large\bfseries\color{mainblue}}{\thesubsection}{0.5em}{}
\newtcolorbox{infobox}[1]{breakable,colback=lightblue,colframe=mainblue,boxrule=0.7pt,arc=1mm,title=\textbf{#1},fonttitle=\bfseries}
\newtcolorbox{notebox}[1]{breakable,colback=lightgray,colframe=black!50,boxrule=0.55pt,arc=1mm,title=\textbf{#1},fonttitle=\bfseries}
\newtheoremstyle{problemstyle}{0.8em}{0.8em}{\normalfont}{}{\bfseries\color{accent}}{．}{0.6em}{}
\theoremstyle{problemstyle}
\newtheorem{exercise}{题目}[section]
\newenvironment{solution}{\par\noindent\textbf{解：}}{\hfill$\square$\par}
\newenvironment{testsolution}{\par\noindent\textbf{参考解答：}}{\hfill$\square$\par}
\newcommand{\answerspace}{\vspace{2.1cm}}
\newcommand{\biganswerspace}{\vspace{3cm}}
\newcommand{\fl}[1]{\left\lfloor #1\right\rfloor}
\newcommand{\fr}[1]{\left\{#1\right\}}

\pagestyle{fancy}\fancyhf{}
\fancyhead[L]{\small 取整与小数部分周测 A 卷}\fancyhead[C]{\small 基础与标准模型}\fancyhead[R]{\small 刘文博}
\fancyfoot[C]{\small 第 \thepage 页 / 共 \pageref{LastPage} 页}
\renewcommand{\headrulewidth}{0.55pt}\renewcommand{\footrulewidth}{0.35pt}
\begin{document}
\thispagestyle{empty}
\begin{center}{\fontsize{22}{28}\selectfont\bfseries 取整函数与小数部分周测 A 卷\par}\vspace{0.35cm}{\large 初中数学联赛专题基础检测\par}\end{center}
\begin{infobox}{考试说明}
\begin{itemize}\item 测试时间：60 分钟\item 满分：100 分\item 独立完成，注意负数取整、整数条件和区间端点。\end{itemize}
\end{infobox}
\section{试题}
\begin{enumerate}
\item（10 分）计算 $[-3.6]$、$\{-3.6\}$ 和 $[\sqrt7]+[-\sqrt7]$。\answerspace
\item（12 分）已知 $[x]=2$，求 $[3x]$ 的所有可能值。\answerspace
\item（12 分）解方程 $x+[x]=8$。\answerspace
\item（12 分）解方程 $[x]+\{2x\}=3$。\answerspace
\item（12 分）解不等式 $-1\le[2x-3]<4$。\answerspace
\item（12 分）解不等式 $[x]>\{x\}$。\answerspace
\item（14 分）证明 $[x]+[y]\le[x+y]\le[x]+[y]+1$。\biganswerspace
\item（16 分）证明对任意实数 $x$，
\[
[x]+[x+\tfrac13]+[x+\tfrac23]=[3x].
\]\biganswerspace
\end{enumerate}
\newpage\section{详细解答}
\begin{enumerate}
\item\begin{testsolution}$[-3.6]=-4$，$\{-3.6\}=0.4$。因 $2<\sqrt7<3$，后两项之和为 $2+(-3)=-1$。\end{testsolution}
\item\begin{testsolution}设 $x=2+t$，$0\le t<1$，则 $[3x]=6+[3t]$，所以可能为 $6,7,8$。\end{testsolution}
\item\begin{testsolution}设 $[x]=n$，则 $x=8-n$ 是整数，故 $[x]=x$，即 $n=8-n$。得 $n=4$、$x=4$。\end{testsolution}
\item\begin{testsolution}设 $x=n+t$。因 $n+\{2t\}=3$ 且第二项既是 $[0,1)$ 内的数，又等于整数 $3-n$，故它为 $0$，$n=3$。于是 $t=0$ 或 $\frac12$，所以 $x=3$ 或 $\frac72$。\end{testsolution}
\item\begin{testsolution}$[2x-3]\ge-1$ 等价于 $2x-3\ge-1$；$[2x-3]<4$ 等价于 $2x-3<4$。故 $x\in[1,\frac72)$。\end{testsolution}
\item\begin{testsolution}设 $x=n+t$，条件为 $n>t$。若 $n\le0$ 无解；若 $n\ge1$，因 $t<1\le n$，全部成立。解集为 $[1,+\infty)$。\end{testsolution}
\item\begin{testsolution}写 $x=[x]+\{x\}$、$y=[y]+\{y\}$，则
\[
[x+y]=[x]+[y]+[\{x\}+\{y\}].
\]
最后一项只能为 $0$ 或 $1$，结论成立。\end{testsolution}
\item\begin{testsolution}设 $x=n+t$。按 $t\in[0,\frac13)$、$[\frac13,\frac23)$、$[\frac23,1)$ 分段，左边依次为 $3n,3n+1,3n+2$，恰与 $[3x]$ 相同。\end{testsolution}
\end{enumerate}
\end{document}
